\documentclass[../../../include/open-logic-section]{subfiles}

\begin{document}

\olfileid[tr]{sth}{replacement}{refproofs}
\olsection{Ek: Yerine Koyma Çevresindeki Sonuçlar}

Bu kesimde Yansıma'yı $\ZF$ içinde ispatlayacağız. Ayrıca Yansıma'nın Yerine
Koyma'ya denk olduğu bir anlamı ispatlayacağız. Son olarak bütün bunların
Yansıma/Yerine Koyma'nın gücüyle ilgili ilginç bir sonucunu ispatlayacağız.
\emph{Uyarı: Bu, bu ders kitabındaki matematiğin açık ara en ileri
kısmıdır.}

Kısa olması için değişken ya da nesne dizilerini ele alırken \emph{üst çizgi}
gösterimini kullanan bir lemayla başlayacağız. Buna göre
``$\overline{a}_k$'', bağlamın belirlediği $n$ için ``$a_{k_1}$, \dots,
$a_{k_n}$'' ifadesini kısaltır.

\begin{lem}\ollabel{lemreflection}
Her $1 \leq i \leq k$ için $\phi_i(\overline{v}_i, x)$ bir formül olsun. O
zaman her $\alpha$ için, bütün $\overline{a}_1, \ldots,
\overline{a}_k \in V_\beta$ ve her $1 \leq i \leq k$ bakımından aşağıdaki
koşulu sağlayan bazı $\beta > \alpha$ vardır:
\[
	\exists x\phi_i(\overline{a}_i, x) \rightarrow (\exists x \in V_\beta) \phi_i(\overline{a}_i, x)
\]
\end{lem}

\begin{proof}
Bir $\mu$ terimini şöyle tanımlarız: $\mu(\overline{a}_1, \ldots,
\overline{a}_k)$, $1 \leq i \leq k$ için aşağıdaki koşulların tümünü
sağlayan en küçük $V$ aşamasıdır:
\[
\exists x\phi_i(\overline{a}_i, x) \rightarrow (\exists x \in V) \phi_i(\overline{a}_i, x)
\]
% [tr source emendation] Upstream has an unmatched final `)` in this display;
% Turkish removes it, matching the lemma statement and intended conditional.
$\mu(\overline{a}_1, \ldots, \overline{a}_k)$ ifadesinin bütün
$\overline{a}_1, \ldots, \overline{a}_k$ için var olduğunu doğrulamak
kolaydır. Şimdi Yerine Koyma ile özyineleme teoremimizi kullanarak şunları
tanımlayın:
\begin{align*}
	S_0 & = V_{\alpha+1}\\
	S_{n+1} & = S_n \cup \bigcup
	\Setabs{\mu(\overline{a}_1, \ldots, \overline{a}_k)}
	{\overline{a}_1, \ldots, \overline{a}_k \in S_n} \\
	S &= \bigcup_{m < \omega} S_m.
\end{align*}
% [tr source emendation] The indexed union upstream ends in `S_n` although
% its bound variable is `m`; Turkish repairs the evident dummy-index error.
Her $S_n$ ve dolayısıyla $S$'nin kendisi $V_\alpha$'dan sonraki bir aşamadır.
Şimdi $\overline{a}_1$, \dots,~$\overline{a}_k \in S$ sabit olsun; öyleyse
$\overline{a}_1$, \dots, $\overline{a}_k \in S_n$ olacak biçimde bazı
$n < \omega$ vardır. Bazı $1 \leq i \leq k$ sabit olsun ve
$\exists x \phi_i(\overline{a}_i,x)$ olduğunu varsayın. Kuruluş gereği
$(\exists x \in \mu(\overline{a}_1, \ldots,
\overline{a}_k))\phi_i(\overline{a}_i, x)$, dolayısıyla
$(\exists x \in S_{n+1})\phi_i(\overline{a}_i, x)$ ve bundan da
$(\exists x \in S)\phi_i(\overline{a}_i, x)$ elde edilir. Demek ki $S$,
aradığımız $V_\beta$'dır.
\end{proof}
\noindent 
Artık \olref[sth][replacement][ref]{reflectionschema} sonucunu oldukça
doğrudan ispatlayabiliriz:

\begin{proof}[İspat] 
$\alpha$ sabit olsun. Genelliği yitirmeden $\phi$'nin yalnızca $\exists$,
$\lnot$ ve $\land$ bağlaçlarını içerdiğini varsayabiliriz (çünkü bunlar ifade
bakımından yeterlidir). $\psi_1, \ldots, \psi_k$, $\phi$'nin alt formüllerini
karmaşıklıklarına göre saysın; böylece $\psi_k = \phi$ olsun.
\olref{lemreflection} gereği, herhangi bir $\overline{a}_i \in V_\beta$ ve
her $1 \leq i \leq k$ için aşağıdaki koşulu sağlayan bir
$\beta > \alpha$ vardır:
\begin{align}\label{reflectionnicelybehaved}
	\exists x\psi_i(\overline{a}_i, x) \rightarrow 
	(\exists x \in V_\beta) \psi_i(\overline{a}_i, x)\tag{*}
\end{align}
$\psi_i$'nin karmaşıklığı üzerinde tümevarımla, herhangi bir
$\overline{a}_i \in V_\beta$ için
$\psi_i(\overline{a}_i) \leftrightarrow
\psi_i^{V_\beta}(\overline{a}_i)$ olduğunu göstereceğiz. $\psi_i$ atomikse
bu önemsizdir. İki koşullu ayrıca, $\psi_i$ bu özelliği sağlayan alt
formüllerin bir değillemesi ya da evetlemesi olduğunda, $\psi_i$'nin
kendisinin de bu özelliği sağladığını kurar. Dolayısıyla tek ilginç durum
nicelemeyle ilgilidir.
$\overline{a}_i \in V_\beta$ sabit olsun; o zaman:
\begin{align*}
	(\exists x \psi_i(\overline{a}_i, x))^{V_\beta}
	&\text{ ancak ve ancak }
	(\exists x \in V_\beta)\psi_i^{V_\beta}(\overline{a}_i, x)
	&&\text{tanım gereği}\\
	&\text{ ancak ve ancak }
	(\exists x \in V_\beta)\psi_i(\overline{a}_i,  x)
	&&\text{varsayım gereği}\\
	&\text{ ancak ve ancak }
	\exists x \psi_i(\overline{a}_i, x)
	&&\text{\eqref{reflectionnicelybehaved} gereği}
\end{align*}
Bu, tümevarımı tamamlar; $\psi_k = \phi$ olduğundan sonuç çıkar.
\end{proof}

Yansıma'yı $\ZF$ içinde ispatladık. İspatımız özünde
\citet{Montague1961}'i izledi. Şimdi $\Z$ içinde Yansıma'nın Yerine Koyma'yı
gerektirdiğini ispatlamak istiyoruz. İspat, bir sadeleştirmeyle birlikte
\citet{Levy1960}'i izler.

$\Z$ içinde çalıştığımızdan Yansıma'yı yukarıda verilen biçimde sunamayız.
Sonuçta Yansıma'yı ``$V_\alpha$'' gösterimini kullanarak formüle ettik; bu
gösterim $\Z$ içinde tanımlanamaz (bkz. \olref[sth][spine][zf]{sec}). Bunun
yerine Yansıma'nın görünüşte daha zayıf bir formülasyonunu şöyle sunacağız:
% [tr source emendation] Upstream says `weaker formulation of Replacement`;
% the immediately defined principle is Weak-Reflection, so Turkish restores
% `weaker formulation of Reflection`.

\begin{defish}
\emph{Zayıf Yansıma.} Her $\phi$ formülü için $0$, $1$ ve $\phi$'nin bütün
parametrelerinin $S$'nin !!{element}s{}ı olduğu, ayrıca $(\forall \overline{x}
\in S)(\phi \liff \phi^S)$ koşulunu sağlayan geçişli bir $S$ kümesi vardır.
\end{defish}

Bunu Yerine Koyma'yı ispatlamakta kullanmak üzere önce \citet[Theorem 2'nin
ilk kısmı]{Levy1960} kaynağını izleyip iki formülü aynı anda
``yansıtabildiğimizi'' göstereceğiz:

\begin{lem}[$\Z + \text{Zayıf Yansıma}$ içinde.]\ollabel{lem:reflect}
Her $\psi, \chi$ formül çifti için $0$ ile $1$'in (ve formüllerin bütün
parametrelerinin) $S$'nin !!{element}s{}ı olduğu, ayrıca $(\forall \overline{x}
\in S)((\psi \liff \psi^S) \land (\chi \liff \chi^S))$ koşulunu sağlayan
geçişli bir $S$ kümesi vardır.
\end{lem}

\begin{proof}
$\phi$, $(z = 0 \land \psi) \lor (z = 1 \land \chi)$ formülü olsun.

Burada bir kısaltma kullanıyoruz; ``$z = 0$'' ifadesini ``$\forall t\, t
\notin z$'', ``$z =1$'' ifadesini de ``$\forall s(s \in z \liff \forall
t\, t \notin s)$'' olarak açıkça yazmalıyız. Fakat $0, 1 \in S$ ve $S$
geçişli olduğundan bu formüller $S$ için \emph{mutlaktır}; yani
niceleyicilerini $S$ ile sınırlasak da aynı nesneye uygulanırlar.\footnote{Daha
biçimsel olarak, $\xi$ bu formüllerden biri olmak üzere $\xi(z) \liff
\xi^S(z)$ olur.}

Zayıf Yansıma gereği aşağıdaki koşulu sağlayan uygun bir $S$ vardır:
\begin{align*}
	(\forall z, \overline{x} \in S)(&\phi \liff \phi^S)\\
	\text{yani }(\forall z, \overline{x} \in S)(&((z = 0 \land \psi) \lor (z = 1 \land \chi)) \liff {}\\
	&\phantom{(}((z = 0 \land \psi) \lor (z = 1 \land \chi))^S)\\
	\text{yani }(\forall z, \overline{x} \in S)(&((z = 0 \land \psi) \lor (z = 1 \land \chi))\liff {}\\
	&\phantom{(}((z = 0 \land \psi^S) \lor (z = 1 \land \chi^S)))\\
	\text{yani }(\forall \overline{x} \in S)(&(\psi \liff \psi^S) \land (\chi \liff \chi^S))
\end{align*}
``$z = 0$'' ile ``$z=1$'' ifadeleri $S$ için mutlak olduğundan ikinci iddia
üçüncüyü; $0 \neq 1$ olduğundan da dördüncü iddia çıkar.
\end{proof}\noindent
Şimdi \citet[Theorem 6]{Levy1960} kaynağını izleyip sadeleştirerek Yerine
Koyma'yı elde edebiliriz:

\begin{thm}[$\Z$ + Zayıf Yansıma içinde]\label{thm:replacement} 
Her $\phi(v,w)$ formülü ve her $A$ için, $(\forall x \in
A)\lexists![y][\phi(x,y)]$ ise $\Setabs{y}{(\exists x \in A)\phi(x,y)}$
vardır.
\end{thm}
% [tr source emendation] Upstream omits the closing parenthesis in the
% displayed set-builder condition; Turkish restores the well-formed formula.

\begin{proof}
$(\forall x \in A)\lexists![y][\phi(x,y)]$ koşulunu sağlayan $A$ sabit olsun
ve şu formülleri tanımlayın:
\begin{align*}
	\psi &\text{ şudur: } (\phi(x, z) \land A = A)\\
	\chi &\text{ şudur: } \lexists[y][\phi(x, y)]
\end{align*}
\olref{lem:reflect} kullanıldığında, $A$, $\psi$'nin bir parametresi
olduğundan $0, 1, A \in S$ (diğer bütün parametrelerle birlikte) olacak ve
şu koşulu sağlayacak biçimde geçişli bir $S$ vardır:
\[
	(\forall x,z \in S)((\psi \liff \psi^S) \land (\chi \liff \chi^S))
\]
Özellikle:
\begin{align*}
	(\forall  x, z \in S)(&\phi(x, z) \liff \phi^S(x, z))\\
	(\forall x \in S)(&\exists y\phi(x, y) \liff (\exists y \in S)\phi^S(x, y)) 
\end{align*}
Bunları birleştirip $A \in S$ ve $S$ geçişli olduğu için $A \subseteq S$
olduğunu gözlemleyince:
\begin{align*}
	(\forall x \in A)(&\exists y\phi(x, y) \liff (\exists y \in S)\phi(x, y))
\end{align*}
elde edilir. $(\forall x \in A)\lexists![y][\phi(x, y)]$ olduğundan şimdi
$(\forall x \in A)(\lexists![y \in S])\phi(x, y)$ olur. Ayırma da
$\Setabs{y \in S}{(\exists x \in A) \phi(x, y)} =
\Setabs{y}{(\exists x \in A) \phi(x, y)}$ sonucunu verir.
\end{proof}

\end{document}
